% Passes 5580--5585: structural identification of the surviving Reye/tomotope
% incidence family as a PSL_2(q) permutation-graph design, with exact Gram/frame
% consequences and a separately firewalled binary-rank conjecture.
\subsection{The Reye/tomotope family as a projectivity-graph incidence design}
\label{sec:pass5580-reye-psl2}

The incidence family isolated in Passes~5488--5495 admits a structural
identification which explains its parameters without an integer coincidence.
For odd prime $q$, write
\[
 X(x)=\begin{pmatrix}x_0&x_2\\-x_3&x_1\end{pmatrix},
 \qquad
 Q(x)=x_0x_1+x_2x_3.
\]
Then $\det X(x)=Q(x)$, so the hyperbolic quadric $Q^+(3,q)$ is the
projectivized rank-one locus.  Writing $X=uv^T$ gives the Segre grid
\[
 \boxed{Q^+(3,q)\cong\mathbf P^1(q)\times\mathbf P^1(q)}
\]
with $(q+1)^2$ points.

For a nonsingular point $p=(a,b,c,d)$ define
\[
 T_p=\begin{pmatrix}c&-a\\-b&-d\end{pmatrix}.
\]
If
\[
 s(u,v)=(u_0v_0,u_1v_1,u_0v_1,-u_1v_0)
\]
is the corresponding singular point, direct expansion of the symplectic form
$B$ gives
\[
 \boxed{B(p,s(u,v))=0\iff u\sim T_pv.}
\]
Moreover $\det T_p=-Q(p)$.  Hence a fixed nonsingular quadratic-value class is
one of the two determinant-square cosets of $\PGL_2(q)$, and after translating
by one fixed projectivity if necessary the incidence structure is exactly the
permutation-graph design of
\[
 G=\PSL_2(q)\curvearrowright\mathbf P^1(q).
\]
Its rows are $g\in G$, its columns are grid cells $(x,y)$, and
\[
 M_{g,(x,y)}=1\iff y=g(x).
\]
Thus
\[
 |G|=\frac{q(q^2-1)}2,
 \qquad
 |\text{columns}|=(q+1)^2,
\]
with row weight $q+1$ and column weight $q(q-1)/2$.  In particular the
$q=3,5,7$ members are precisely
\[
 12_4\,16_3,\qquad 60_6\,36_{10},\qquad 168_8\,64_{21}.
\]
The first is the Reye configuration/tomotope medial layer already identified
in Pass~5490; the new statement is the projectivity-graph mechanism for the
whole tested family.

\paragraph{Rook-complement Gram theorem.}
Let $A$ be the adjacency matrix of the complement of the
$(q+1)\times(q+1)$ rook graph.  Two distinct cells can lie on one projectivity
graph precisely when both coordinates differ; $2$-transitivity then leaves a
two-point stabilizer of size $(q-1)/2$.  Therefore
\[
 \boxed{
 M^TM=\frac{q-1}{2}(qI+A).
 }
\]
Since $A$ is the tensor graph $K_{q+1}\times K_{q+1}$ with spectrum
\[
 q^2{}^{\,1}\oplus1^{q^2}\oplus(-q)^{2q},
\]
one obtains
\[
 \boxed{
 \operatorname{spec}(M^TM)=
 \left(\frac{q(q^2-1)}2\right)^1
 \oplus
 \left(\frac{q^2-1}{2}\right)^{q^2}
 \oplus0^{2q},
 }
\]
and hence
\[
 \boxed{\operatorname{rank}_{\mathbf Q}M=q^2+1.}
\]
The executable certificate independently reaches this rank modulo $1000003$
for $q=3,5,7,11,13$.

\paragraph{Centered tight frame.}
Every row has mean $1/(q+1)$ and centered squared norm $q$.  Distinct
projectivities agree at $t\in\{0,1,2\}$ points, so after unit normalization the
centered inner products lie in
\[
 \boxed{\left\{-\frac1q,0,\frac1q\right\}.}
\]
Removing the constant eigenline leaves one nonzero Gram eigenvalue, so the
$|\PSL_2(q)|$ centered rows form a unit-norm tight frame in $\mathbf R^{q^2}$
with frame bound
\[
 \boxed{\frac{q^2-1}{2q}.}
\]
At $q=3$ only $-1/3$ and $0$ occur, giving a two-distance specialization.

\paragraph{The order-$576$ Reye symmetry becomes explicit.}
At $q=3$, $\PSL_2(3)\cong A_4$ on four letters.  The Reye incidence is
\[
 g\sim(x,y)\iff y=g(x),
 \qquad g\in A_4,\quad(x,y)\in[4]^2.
\]
Every pair $(a,b)\in S_4^2$ with equal parity acts by
\[
 (x,y)\mapsto(ax,by),
 \qquad g\mapsto bga^{-1},
\]
giving $24^2/2=288$ automorphisms.  Adjoining
$(x,y)\leftrightarrow(y,x)$ and $g\leftrightarrow g^{-1}$ doubles this to
\[
 \boxed{576}
\]
with abstract shape
\[
 \boxed{
 \{(a,b)\in S_4^2:\operatorname{sgn}(a)=\operatorname{sgn}(b)\}\rtimes C_2
 \cong ((A_4\times A_4):C_2):C_2.
 }
\]
This is the same abstract group shape that Pass~5516 obtained by GAP for the
Klein-$V_4$ Latin-square autoparatopy/13-cover image.  This observation explains
why the Reye object naturally carries that $576$-group; it does \emph{not}
identify the underlying acted-on objects.

\paragraph{Characteristic-two firewall.}
The measured binary ranks are
\[
 8,18,32,72,98
 \qquad(q=3,5,7,11,13),
\]
which equal $(q+1)^2/2$ in every tested case.  We record
\[
 \operatorname{rank}_2M\stackrel{?}=\frac{(q+1)^2}{2}
\]
as a conjecture, not an all-$q$ theorem.  The characteristic-two nonsplitting
warning of Passes~5356--5361 applies directly; a proof must analyze the image
algebra of the modular projective-line permutation representation rather than
reduce the characteristic-zero Gram decomposition.

\paragraph{Evidence boundary.}
The coordinate incidence identity, rook-complement Gram factorization, and
$q=3$ order-$576$ subgroup are exact.  The committed verifier replays the
construction for the five odd primes $3,5,7,11,13$.  The prime-power extension,
the binary-rank law beyond these anchors, and any polytope realization for
$q>3$ remain open.  No physical claim follows from this finite incidence
identification.
